\chapter*{Technical Abstract}
\addcontentsline{toc}{chapter}{Technical Abstract}

This project investigates whether three-dimensional active acoustic localisation can be performed using a structured neuromorphic model inspired by bat echolocation, rather than an end-to-end black-box regressor. The task is to estimate target distance, azimuth, and elevation from the echo produced by a known emitted call. This is a relevant problem for edge sensing systems because localisation must often be performed under constraints on computation, latency, memory, and power. The biological auditory system provides a useful design reference because echolocating bats solve a closely related active sensing problem using sparse neural processing and specialised parallel pathways.

The model developed in this project is a full simulated echolocation pipeline. A frequency-modulated call is emitted into a controlled acoustic environment, and the returning binaural echo is generated with distance-dependent delay, propagation attenuation, interaural timing and level differences, and an elevation-dependent comb-filter spectral cue. The received waveform is processed by an IIR cochlear filterbank and a dynamic leaky integrate-and-fire spike encoder, producing a tonotopic event-based representation. This representation is then passed through three parallel pathway models. Distance is estimated using echo delay, corollary discharge, onset extraction, and delay coincidence. Azimuth is represented using both interaural time-difference coincidence and interaural level-difference opponent comparison. Elevation is represented using a dorsal-cochlear-nucleus-like spectral-template model that compares the observed selected-ear spectrum with expected comb-filter transfer functions.

The pathway outputs are represented as one-dimensional neural populations. Continuous Attractor Neural Networks are tested as superior-colliculus-like population readouts for distance, azimuth, and elevation. These attractor readouts are used as stabilising neural coordinate representations rather than as general-purpose optimisers. Their effect is pathway-dependent: they slightly improve the constrained distance readout, are not always beneficial for azimuth, and are most useful for elevation when combined with an inverse-sigmoid calibration that corrects a stable monotonic readout distortion. The final model therefore preserves both raw pathway populations and CANN readouts.

A small trainable spiking neural network is then added as a final fusion stage. It receives raw distance, azimuth, and elevation populations, CANN coordinate readouts, and confidence features such as population activation and spike-count statistics. The output is encoded as distance and sine/cosine angle components. Two trainable variants are compared: a direct readout from the cached pathway features and a residual readout that learns a correction to the fixed pathway estimate. Additional SNN baselines are trained from less structured inputs, including raw binaural waveforms and cochlear spike rasters, to test whether the hand-designed auditory pathways provide useful feature extraction.

The results are reported using coordinate-specific mean absolute errors, Euclidean position error, and a dimensionless combined error. The combined error is formed from normalised distance error and angular sine/cosine error, and is used only to compare readout variants within this project.

In the constrained operating region of 0.25--5~m and $\pm45^\circ$ azimuth/elevation, the fixed pathway model already produces functional localisation. In the clean case, the raw pathway centre-of-mass readout gives distance, azimuth, and elevation mean absolute errors of 0.0365~m, 4.580$^\circ$, and 2.986$^\circ$ respectively. Adding the fixed CANN/calibrated readout improves elevation error to 0.910$^\circ$ and reduces Euclidean position error from 0.2999~m to 0.2435~m. The residual SNN readout gives the best clean result, reducing combined error to 0.0223 and Euclidean mean absolute error to 0.1247~m.

Environmental noise exposes the main limitation of the fixed readout. When a 50~dB environmental noise field is added before binaural and elevation filtering, the fixed CANN baseline suffers a large elevation error, with combined error increasing to 0.2626 and Euclidean error to 1.5772~m. The residual SNN substantially corrects this failure, reducing combined error to 0.0548 and Euclidean error to 0.2968~m. When reverberant/delayed echo components are added alongside the same environmental noise, the residual SNN remains stable, with combined error 0.0550 and Euclidean error 0.3103~m. These results indicate that the trainable readout is using the wider pathway feature set to correct systematic cue interactions rather than simply replacing the underlying model.

The input-only SNN baselines support the same conclusion. Raw waveform SNNs perform poorly despite receiving the original acoustic signal, while cochlear-raster SNNs perform substantially better, showing that early auditory preprocessing is useful. However, the best overall performance is obtained when the SNN receives the full set of biologically motivated pathway features. The project therefore demonstrates a hybrid design principle: fixed biomimetic pathways provide interpretable and useful feature extraction, while a small trainable SNN provides contextual integration and residual correction.

The current system remains a constrained single-object, single-pulse prototype. It uses simplified acoustic cues, functional rather than biophysical neural models, and hand-tuned pathway parameters. Nevertheless, the model establishes a complete foundation for extending neuromorphic echolocation toward continuous multi-pulse tracking, active feedback control, velocity estimation, and more realistic multi-object acoustic scenes.
